\section{Uniform character seeds and the common abelian CM field}
\label{sec:common-arithmetic-frame-v50}

This section constructs the character-and-field input to the ordinary
automorphy argument.  Its principal endpoint is the full original
\texttt{seed\_cl}: one modulus is chosen before the later abelian CM
field and auxiliary prime, and every compatible later choice receives
the entire original \texttt{SeedClause}.  We also construct the actual
fixed field in \texttt{field\_ex}, the auxiliary prime in
\texttt{auxw\_ex}, and install the seed on an already supplied common
frame without changing its marked element.  The seed here is a triple
of characters.  Constructing the rank-two and rank-three automorphic
representations from that triple remains a separate task.

\subsection{Objects, hypotheses, and the order of choices}

Fix the supplied algebraic closure \(\Omega/\mathbb Q\), put
\(G=\operatorname{Gal}(\Omega/\mathbb Q)\), and let \(\ell\) be an
odd prime.  The coefficient ring \(O=c.O\subset\overline{\mathbb Q}_\ell\)
is the original finite \(\mathbb Z_\ell\)-algebra which is a DVR.
Write \(k=O/\mathfrak m_O\), \(q=|k|\), and
\[
 \epsilon:G\longrightarrow\mathbb Z_\ell^\times,\qquad
 \omega:G\longrightarrow k^\times,
 \qquad \alpha=\overline\Psi:G\longrightarrow k^\times .
 \tag{CAF.1}
\]
Here \(\omega\) is the reduction of the actual cyclotomic character,
and \(\Psi:G\to O^\times\) and
\(\eta:G\to\{\pm1\}\) are the specified characters.  A complex
conjugation \(\sigma\in G\) and a valuation \(B\) over \(\ell\)
are fixed.  For a positive integer \(n\), the source's
\(\operatorname{KillsMu}(n,\chi)\) means
\[
 s|_{\mu_n}=1\quad\Longrightarrow\quad \chi(s)=1.
 \tag{CAF.2}
\]
It is a factorization through the actual cyclotomic extension inside
\(\Omega\), rather than an assertion about an unrelated character of
the same order.  Decomposition and inertia subgroups at a valuation
\(A\) are denoted by \(D_A,I_A\).

The full original \(\operatorname{Gen}(\ell,\eta,\Psi,\sigma,B,M,H_b)\)
is an allowed starting input.  For the seed construction only the
following clauses of that record are needed:
\begin{enumerate}
\item \(k\) is finite, \(M>0\), \(\ell\mid M\), \(\ell\ne2\),
      \(\sigma\) is a complex conjugation, and \(B\) lies over \(\ell\).
\item There are \(D>0\) and \(\delta\in\Omega\) such that
      \(\delta^2=-D\) and \(s\delta=\eta(s)\delta\) for every \(s\in G\).
\item For one integer \(a\ge0\), both \(\eta\) and \(\alpha\)
      satisfy \(\operatorname{KillsMu}(M^a,-)\).
\item \(\eta|_{I_A}=1\) for every \(A\) over \(\ell\),
      \(\alpha(\sigma)=1\), and there is one \(j\ge0\) such that
      \(\alpha|_{I_A}=\omega^j|_{I_A}\) for every such \(A\).
\item \(\ker\omega\cap\ker\alpha\not\subseteq\ker\eta\).
\end{enumerate}
Call these explicit clauses \(\mathcal S\).  Every use of
\(\mathcal S\) below is therefore available from the original
\texttt{Gen} without an extra hypothesis.  Conversely, proving the
seed from these fewer clauses does not assert the other
\texttt{Gen} fields.  In particular, \(H_b\), the local degree bounds,
and the strict cohomology classes are not constructed by the seed
argument.  Claims CAF9 and CAF10 verify \(\mathcal S\) directly for
the original \texttt{Dat} and \texttt{HostData} applications.

For an abelian CM field \(L\subset\Omega\), let
\(U_L=G_{L^+}\), the source's \texttt{upl}, and
\(G_L=\operatorname{Gal}(\Omega/L)\).  Put
\[
 H_\ell(w)=\{s\in G:\ s|_{\mu_w}\text{ is an }\ell
                  \text{-th power in }(\mathbb Z/w\mathbb Z)^\times\}.
 \tag{CAF.3}
\]
This equals the literal \texttt{Hth}: the source additionally writes
the same condition for \(s^{-1}\), which is automatic because the
\(\ell\)-th powers form a subgroup.

A \emph{witness prime} is exactly the original
\texttt{FrobeniusWitnessPrime}: a prime \(w\nmid M\) with
\(w\equiv1\pmod\ell\), such that for every \(A\) above \(w\),
\[
          D_A\subseteq\ker\eta\cap\ker\alpha\cap\ker\omega.
 \tag{CAF.4}
\]
All assertions below use the same \(\Omega,O,\eta,\Psi,\sigma\).
The embedding \(\iota:\overline{\mathbb Q}_\ell\simeq\mathbb C\)
needed for later automorphy is not used or changed here.

\subsection{Local algebra and the two cyclotomic constructions}

% BEGIN CAF1
\paragraph{CAF1. Teichmuller lifts with the original norm congruence.}
\label{claim:CAF1}
There is an injective homomorphism
\(T:k^\times\to\overline{\mathbb Q}_\ell^\times\), landing in
\(O^\times\), such that for every \(u\in O^\times\), every
\(j\ge0\), and every \(s\in G\),
\[
 \left|\epsilon(s)^j
       T\bigl(\overline u\,\omega(s)^{-j}\bigr)-u\right|_\ell<1.
 \tag{CAF.5}
\]
Moreover, \(|x|_\ell\le1\) for every \(x\in O\), and
\(x\in\mathfrak m_O\) if and only if \(|x|_\ell<1\).
These assertions need only that \(\ell\) is prime, including
\(\ell=2\); the parity observation at the end of the proof uses
oddness.

\emph{Proof.}
Integrality over \(\mathbb Z_\ell\) gives the first norm bound: in a
monic integral equation, a root of norm greater than one would make
the leading term strictly larger than every other term, contradicting
the ultrametric inequality.  A unit and its inverse both have norm at
most one, so both have norm one.  If \(\varpi\) is a uniformizer,
write \(\ell=u\varpi^r\) in the DVR.  Since
\(|\ell|_\ell=\ell^{-1}<1\) and \(|u|_\ell=1\),
\(|\varpi|_\ell<1\).  The factorization of any nonzero element as a
unit times a power of \(\varpi\) now proves the asserted equivalence,
including zero separately.

The ring \(O\) is complete: it is finite over the complete DVR
\(\mathbb Z_\ell\), and its maximal-ideal topology is the same as
its \(\ell\)-adic topology.  For \(b\in k^\times\), Hensel's lemma
gives a unique root \(T(b)\in O\) of \(X^{q-1}-1\) reducing to
\(b\), since the derivative is a unit modulo \(\mathfrak m_O\).
Uniqueness shows \(T(bb')=T(b)T(b')\), and reduction proves
injectivity.  The reduction of
\(\epsilon(s)^jT(\overline u\,\omega(s)^{-j})\) is precisely
\(\overline u\).  Its difference from \(u\) belongs to
\(\mathfrak m_O\), proving (CAF.5).  Since \(\ell\) is odd,
\(-1\ne1\) in \(k\), and uniqueness also gives \(T(-1)=-1\).
This proves the complete original \texttt{sc\_norm} and
\texttt{sc\_teich} conclusions.
% END CAF1

% BEGIN CAF2
\paragraph{CAF2. Removing the cyclotomic conductor at \(\ell\).}
\label{claim:CAF2}
Let \(C\) be any group and \(\chi:G\to C\) a homomorphism.
If \(N_0>0\), \(\ell\nmid N_0\),
\(\operatorname{KillsMu}(\ell^rN_0,\chi)\), and
\(\chi|_{I_B}=1\), then
\(\operatorname{KillsMu}(N_0,\chi)\).
If also \(w\equiv1\pmod {N_0}\), then the source's full
\(\operatorname{SplitAt}(N_0,w,\chi)\) holds.

\emph{Proof.}
Inertia at \(\ell\) acts trivially on roots of unity of order prime
to \(\ell\).  Indeed, those roots are the Hensel lifts of the roots
in finite residue-field extensions, and hence lie in the maximal
unramified extension.  In contrast,
\(\mathbb Q_\ell(\mu_{\ell^r})/\mathbb Q_\ell\) is totally
ramified: for \(r\ge1\),
\(\Phi_{\ell^r}(1+X)\) is Eisenstein at \(\ell\), of degree
\(\ell^{r-1}(\ell-1)\).  Inertia therefore surjects onto the full
group of its cyclotomic automorphisms.  The case \(r=0\) is immediate.

Suppose \(s\) fixes \(\mu_{N_0}\).  Choose \(t\in I_B\) with the
same action as \(s\) on \(\mu_{\ell^r}\).  It fixes
\(\mu_{N_0}\), so \(t^{-1}s\) fixes both sets of roots.  Since the
orders are coprime, these roots generate \(\mu_{\ell^rN_0}\).
Thus \(\chi(t^{-1}s)=1\), and \(\chi(t)=1\) gives \(\chi(s)=1\).
No commutativity of \(C\) was used.  Finally, an element acting as
\(\zeta\mapsto\zeta^w\) on \(\mu_{N_0}\) fixes it when
\(w\equiv1\pmod {N_0}\), so \(\chi\) is one on every such
element.  These are exactly the two \texttt{SplitAt} clauses.
The proof preserves the one given valuation \(B\) and supplies the
original \texttt{sc\_strip} for every target group.
% END CAF2

% BEGIN CAF3
\paragraph{CAF3. The actual imaginary quadratic field and its discriminant.}
\label{claim:CAF3}
The quadratic clause of \(\mathcal S\) alone makes the subfield
\(K_0=\mathbb Q(\delta)\) of \(\Omega\) imaginary quadratic with
\(G_{K_0}=\ker\eta\).  With the inertia clause as well,
\[
                 G_{K_0}=\ker\eta,\qquad
                 \ell\nmid\operatorname{disc}(K_0).
 \tag{CAF.6}
\]
More generally, the discriminant conclusion holds for any supplied
quadratic field \(K\) for which every inertia subgroup above
\(\ell\) lies in \(G_K\).  No oddness assumption is needed for
this assertion.

\emph{Proof.}
The polynomial \(X^2+D\) has no rational root because \(D>0\).
It is therefore the minimal polynomial of \(\delta\), and all its
complex roots are nonreal.  An element \(s\in G\) fixes \(K_0\)
exactly when it fixes \(\delta\).  As \(\delta\ne0\) and
\(s\delta=\eta(s)\delta\), this is equivalent to \(\eta(s)=1\).
Every inertia subgroup above \(\ell\) lies in that kernel by
\(\mathcal S\), so \(K_0\) is unramified above \(\ell\).

For completeness, the discriminant consequence uses the ordinary
integral trace pairing.  The completion of the ring of integers at
\(\ell\) is a product of the integer rings of the local extensions.
Trivial inertia makes each of these unramified; its residue extension
is a finite separable field extension.  Reduction of its trace pairing
is the nondegenerate field-trace pairing over the finite residue field.
Its determinant is consequently a unit.  Taking the product over the
places shows that the global discriminant is a unit at \(\ell\).
This is the source's \texttt{sc\_quad} and \texttt{sc\_disc}, on the
actual embedded quadratic field.
% END CAF3

% BEGIN CAF4
\paragraph{CAF4. The prescribed order-\(\ell\) character at \(w\).}
\label{claim:CAF4}
For primes \(\ell,w\) with \(w\equiv1\pmod\ell\), there is a
character \(\theta:G\to\overline{\mathbb Q}_\ell^\times\) such that
\[
 \ker\theta=H_\ell(w),\qquad \theta\ne1,
 \qquad\theta^\ell=1,
 \qquad\operatorname{KillsMu}(w,\theta).
 \tag{CAF.7}
\]
The group \(H_\ell(w)\) is open, contains \([G,G]\), has index
\(\ell\), contains all inertia away from \(w\), and does not
contain inertia at any place over \(w\).  These assertions hold also
for \(\ell=2\); only the further real-place observation below uses
oddness.

\emph{Proof.}
The cyclotomic map \(G\twoheadrightarrow
(\mathbb Z/w\mathbb Z)^\times\) is onto; the target is cyclic of
order \(w-1\).  The subgroup of \(\ell\)-th powers has index
\(\ell\).  Compose this map, its quotient by the power subgroup,
and an isomorphism of that cyclic quotient with
\(\mu_\ell(\overline{\mathbb Q}_\ell)\).  This gives \(\theta\)
and all four identities in (CAF.7).  It also proves openness and
the commutator assertion.  Roots of order \(w\) are unramified away
from \(w\), while total ramification of
\(\mathbb Q_w(\mu_w)\) makes inertia at \(w\) surject onto the
entire cyclotomic group.  Its image in the order-\(\ell\) quotient
is therefore nontrivial at every prolongation of the place.  For odd
\(\ell\), complex conjugation belongs to \(H_\ell(w)\), because
its cyclotomic image \(-1\) is \((-1)^\ell\).
This proves \texttt{hth} and \texttt{sc\_theta}, including the
literal kernel equality needed in the later field choice.
% END CAF4

% BEGIN CAF5
\paragraph{CAF5. A witness prime satisfying every prescribed congruence.}
\label{claim:CAF5}
Given the original \texttt{Gen} and any \(Q>0\), there is a witness
prime \(w\) satisfying \(w\equiv1\pmod Q\).  In fact, the
weaker clauses \(\mathcal S\) suffice.

\emph{Proof.}
Choose \(a\) in \(\mathcal S\) and put \(R=M^{a+1}Q\).  There
are arbitrarily large primes congruent to one modulo \(R\).
Only this residue class is needed, so the elementary cyclotomic proof
suffices: for \(R>1\), take a sufficiently large integer \(A\)
divisible by \(R\)
and by any proposed finite list of primes.  A prime divisor of
\(\Phi_R(A)>1\) divides neither \(A\) nor \(R\), since
\(\Phi_R(0)=1\).  The order of \(A\) modulo that prime is exactly
\(R\): the cyclotomic factors of \(X^R-1\) have disjoint roots
in characteristic prime to \(R\).  Hence the prime is one modulo
\(R\) and is outside the proposed list.  The case \(R=1\) is
Euclid's argument.  Select such a prime \(w>R\).

Then \(w\nmid M\), \(w\equiv1\pmod\ell\), and
\(w\equiv1\pmod Q\).  For \(A\) above \(w\), decomposition
acts on \(\mu_n\), \(n\mid R\), through powers of the Frobenius
\(\zeta\mapsto\zeta^w\).  These roots are unramified there, and
\(w\equiv1\pmod n\), so the entire decomposition subgroup fixes
them.  In particular it fixes \(\mu_{M^a}\) and \(\mu_\ell\).
The two \texttt{KillsMu} conditions and the definition of \(\omega\)
give (CAF.4).  This proves \texttt{auxw\_ex} without a stronger
distribution theorem for primes or a field-dependent modulus.
% END CAF5

\subsection{One modulus before every later field and prime}

% BEGIN CAF6
\paragraph{CAF6. Fixed core characters and a prime-to-\(\ell\) modulus.}
\label{claim:CAF6}
Choose \(a,j,\delta\) from \(\mathcal S\), put
\[
 M^a\ell=\ell^rN_0,\quad \ell\nmid N_0,\quad N_0>0,
 \qquad \phi=\alpha\omega^{-j},\quad
 \chi'=T\phi,\quad \eta'=\eta\text{ viewed in }
                  \overline{\mathbb Q}_\ell^\times.
 \tag{CAF.8}
\]
All of \(K_0,N_0,j,\chi',\eta'\) are fixed before \(L,w\).
They satisfy
\[
 \begin{gathered}
 \operatorname{KillsMu}(N_0,\chi'),\qquad
 \operatorname{KillsMu}(N_0,\eta'),\qquad
 \ell\nmid\operatorname{ord}(\chi'),\qquad \chi'\ne\eta',\\
 \ker\eta'=G_{K_0},\qquad \chi'(\sigma)=(-1)^j,\qquad
 |(\epsilon^j\chi')(s)-\Psi(s)|_\ell<1\quad(s\in G).
 \end{gathered}
 \tag{CAF.9}
\]

\emph{Proof.}
The characters \(\eta,\alpha\) factor through \(\mu_{M^a}\),
and \(\omega\) factors through \(\mu_\ell\).  Thus \(\chi'\)
and \(\eta'\) satisfy \(\operatorname{KillsMu}(M^a\ell,-)\).
At every inertia group over \(\ell\),
\(\phi=\alpha\omega^{-j}=1\), and \(\eta'=1\).
CAF2 strips the factor \(\ell^r\) and proves the two factorizations
through \(\mu_{N_0}\).

The image of \(\chi'\) lies in the Teichmuller group of order
\(q-1\).  Its order divides \(q-1\), which is prime to \(\ell\).
The natural embedding of \(\{\pm1\}\) in the characteristic-zero
field is injective, so \(\ker\eta'=\ker\eta=G_{K_0}\).
At \(\sigma\), \(\alpha=1\) and \(\omega=-1\), giving
\(\chi'(\sigma)=T((-1)^{-j})=(-1)^j\).  Equation (CAF.5), with
\(u=\Psi(s)\), gives the last assertion of (CAF.9) for every
\(s\), with the original integral embedding of \(O\).

Finally, choose \(s\in\ker\omega\cap\ker\alpha\) outside
\(\ker\eta\), using \(\mathcal S\).  At this element
\(\chi'(s)=T(1)=1\) and \(\eta'(s)=-1\).  Hence
\(\chi'\ne\eta'\).  No later field or representation has been
chosen in this proof.
% END CAF6

% BEGIN CAF7
\paragraph{CAF7. The complete original uniform \texttt{SeedClause}.}
\label{claim:CAF7}
Set \(Q=N_0\) from CAF6.  For every later abelian CM field
\(L\subset\Omega\) and every later witness prime \(w\), assume
\[
 G_L=U_L\cap\ker\eta,\qquad w\equiv1\pmod Q,
 \qquad U_L\subseteq H_\ell(w),\qquad
 \exists s_0\in G_L:\alpha(s_0)^2\ne1.
 \tag{CAF.10}
\]
Then the full original
\(\operatorname{SeedClause}(c,\eta,\Psi,\sigma,L)\) holds.
Consequently this proves \texttt{seed\_cl} for every original
\texttt{Gen}, with its entire \(\forall L\,\forall w\) scope.

\emph{Proof.}
Use \(K_0,N_0,j,\chi',\eta'\) already fixed in CAF6, and choose
\(\theta\) from CAF4 for this \(w\).  The numerical conditions in
\texttt{SlopeTestRemainingChoices} are all satisfied: \(\ell\ne2\),
\(N_0>0\), \(\ell\nmid N_0\), \(w\) is prime, and
\(w\nmid N_0\).  For the last assertion, \(N_0\mid M^a\ell\),
whereas \(w\nmid M\) and \(w\ne\ell\); the latter follows from
\(w\equiv1\pmod\ell\).  The remaining three auxiliary conditions
\(\theta\ne1\), \(\theta^\ell=1\), and
\(\operatorname{KillsMu}(w,\theta)\) are exactly (CAF.7).

We check every field condition.  The given \(\sigma\) is a complex
conjugation.  CAF3 gives \([K_0:\mathbb Q]=2\) and imaginary
quadratic \(K_0\).  Abelianity of \(L\) is supplied in (CAF.10).
The equality there gives
\(G_L\subseteq\ker\eta=G_{K_0}\), which is the original
\texttt{KL} inclusion, and CAF6 gives
\(\ker\eta'=G_{K_0}\).  Since \(w\equiv1\pmod {N_0}\),
CAF2 and CAF6 give both
\(\operatorname{SplitAt}(N_0,w,\eta')\) and
\(\operatorname{SplitAt}(N_0,w,\chi')\).  Furthermore,
\[
        G_L\subseteq U_L\subseteq H_\ell(w)=\ker\theta,
\]
which is the original \texttt{thetaL} condition, and the parity
condition is \(\chi'(\sigma)=(-1)^j\) from CAF6.

For the remaining square condition, let
\(z=(\epsilon^j\chi')(s_0)\) and \(u=\Psi(s_0)\).
Equation (CAF.9) gives \(|z-u|_\ell<1\).  Since \(u\) is a
unit, both \(|u|_\ell=1\) and \(|z|_\ell=1\).
If \(|z^2-1|_\ell<1\), then
\[
 u^2-1=(u-z)(u+z)+(z^2-1)
\]
has norm less than one.  CAF1 would give
\(\alpha(s_0)^2=1\), contrary to (CAF.10).  Therefore
\(\neg\operatorname{Cong}(z^2,1)\), precisely the source's
\texttt{sq} clause.  This proves every field of
\texttt{FieldPair\allowbreak{}Character\allowbreak{}TripleConditions}.

The four additional \texttt{SeedClause} requirements are
\(\chi'\ne\eta'\), \(\ell\nmid\operatorname{ord}(\chi')\),
the literal equality of values of \(\eta'\) with the integer-valued
\(\eta\), and the all-\(s\) norm congruence.  They are (CAF.8)--(CAF.9).
The final discriminant requirement \(\ell\nmid\operatorname{disc}K_0\)
is CAF3.  Nothing in these choices depends on a later automorphic
representation, coefficient extension, or universal-ring point.
The proof of the square implication also proves the original
\texttt{sc\_cong} on its full stated inputs.
% END CAF7

\subsection{The actual CM field, and the two arithmetic applications}

% BEGIN CAF8
\paragraph{CAF8. An explicit common abelian CM field from the prescribed subgroup.}
\label{claim:CAF8}
Let \(U\subseteq G\) be open, contain \([G,G]\), and contain
every complex conjugation.  Under the quadratic clause of
\(\mathcal S\), define the actual subfield
\[
                   L=\Omega^{U\cap\ker\eta}.
 \tag{CAF.11}
\]
Then \(L/\mathbb Q\) is finite abelian and CM, and literally
\[
             G_L=U\cap\ker\eta,\qquad U_L=U.
 \tag{CAF.12}
\]
Thus the original \texttt{field\_ex} is proved for every original
\texttt{Gen}, every \(m\), and every supplied
\(\operatorname{UFacts}(\ell,\eta,\Psi,M,m,U)\).

\emph{Proof.}
A subgroup containing the commutator is normal, and its quotient is
abelian.  Both \(U\) and \(\ker\eta=G_{K_0}\) are open, so
infinite Galois correspondence identifies their fixed fields with
finite Galois extensions in the given \(\Omega\).  Put
\(F=\Omega^U\).  Every complex embedding of \(F\) extends to an
embedding of \(\Omega\) in \(\mathbb C\).  The corresponding
complex conjugation lies in \(U\), so fixes \(F\).  Hence every
such embedding has real image and \(F\) is totally real.

The fixed field of an intersection of these closed subgroups is the
compositum, so \(L=FK_0\) inside \(\Omega\).  Since the imaginary
quadratic field \(K_0\) has no nontrivial real subfield,
\(F\cap K_0=\mathbb Q\).  Thus \([L:F]=2\), and every complex
embedding of \(L\) contains a nonreal root of \(X^2+D\).
The extension \(L/F\) is totally imaginary quadratic, so it is CM
with maximal real subfield exactly \(F\).  Consequently
\(U_L=G_F=U\), and the first equality in (CAF.12) follows directly
from the definition (CAF.11).  Both groups and their embeddings are
the original subgroups of \(G\); there is no replacement by an
abstract isomorphic CM field.  In \texttt{UFacts}, precisely
\texttt{isOpen}, \texttt{ab}, and \texttt{real} provide the three
subgroup hypotheses used here.  Its remaining local clauses are
retained for subsequent frame construction.
% END CAF8

% BEGIN CAF9
\paragraph{CAF9. The seed hypotheses for the original \texttt{Dat}.}
\label{claim:CAF9}
For the original \(D:\operatorname{Dat}(\Omega,K,c,e,\psi)\),
use its characters \(\eta=D.\eta\), \(\Psi=\epsilon^e\psi\),
its conjugation \(\sigma=D.\mathrm{cc}\), and any supplied \(B\)
over \(\ell\).  Set
\[
 M=\ell N|\operatorname{disc}K|,
 \qquad Q=|\operatorname{disc}K|\,N_{\ell'},
 \tag{CAF.13}
\]
where \(N_{\ell'}\) is the prime-to-\(\ell\) part of \(N\).
Then \(\mathcal S\) holds, and CAF7 supplies the original
\texttt{SeedClause} for every later \(L,w\) satisfying (CAF.10)
with this \(Q\).  The quadratic field in that clause may be the
original embedded \(K\).

\emph{Proof.}
The original \(GD\) conditions give imaginary quadratic \(K\),
\(\ell\ne2\), even positive \(N\),
\(\ell\nmid\operatorname{disc}K\), prime-to-\(\ell\) order of
\(\psi\), \((\operatorname{comp}_\ell\psi)^2=1\), and
\(\psi(-1)=(-1)^e\).  The ring \(O\) is finite over
\(\mathbb Z_\ell\), so its residue field is finite.
Write \(\Delta=\operatorname{disc}K<0\).  The discriminant
determinant of an integral basis is an element \(\delta\in K\)
with \(\delta^2=\Delta=-|\Delta|\).  It generates the quadratic
field, and its two conjugates are \(\delta,-\delta\).  Since
\(\ker\eta=G_K\), \(s\delta=\eta(s)\delta\).

The elementary quadratic conductor theorem places
\(K\) in \(\mathbb Q(\mu_{|\Delta|})\), so
\(\operatorname{KillsMu}(|\Delta|,\eta)\).  One usual explicit
proof uses the primitive quadratic Kronecker character \(\chi_\Delta\):
its Gauss sum \(g=\sum_{a\bmod |\Delta|}\chi_\Delta(a)\zeta^a\)
satisfies \(s_b(g)=\chi_\Delta(b)g\) and \(g^2=\Delta\).
The latter identity follows from finite additive-character
orthogonality, giving \(|g|^2=|\Delta|\), and
\(\overline g=\chi_\Delta(-1)g=-g\).
Thus the same quadratic field generated by \(\delta\) is contained
in that cyclotomic field.  This is the usual elementary quadratic
conductor input; no general potential modularity is used.
Also \(\psi\) factors through \(\mu_N\), and \(\omega\) through
\(\mu_\ell\).  Hence \(\eta\) and
\(\alpha=\omega^e\overline\psi\) both satisfy
\(\operatorname{KillsMu}(M,-)\), so \(a=1\) is allowed.

The character \(\eta\) is unramified at \(\ell\).
At \(\ell\), the residual character of the finite local component
of \(\psi\) factors through \(\mathbb F_\ell^\times\): the
\(\ell\)-Sylow subgroup maps trivially into the prime-to-\(\ell\)
group \(k^\times\).  Therefore it is \(\omega^b\) for one
\(0\le b<\ell-1\).  The same exponent works at every prolongation
of the place, by conjugacy and commutativity of the target.  Taking
\(j=e+b\) proves the inertia clause, and the stated parity gives
\(\alpha(\sigma)=1\).

It remains to verify the kernel noncontainment.  The original
\(GD.\eta\) gives an \(s\) with \(\omega(s)=1\),
\(\chi_P(s)=1\), and \(\eta(s)=-1\), where, since \(\psi\)
has order prime to \(\ell\),
\(\chi_P=\psi(\operatorname{comp}_\ell\psi)^{-1}\).
The local component has order prime to \(\ell\), and therefore
factors through the mod-\(\ell\) cyclotomic character rather than
its pro-\(\ell\) kernel.  It is one at this \(s\).
Consequently \(\psi(s)=1\), and \(\alpha(s)=1\), as required.
CAF6 with \(a=1\) yields
\(Q=(M\ell)_{\ell'}=|\Delta|N_{\ell'}\), exactly (CAF.13).
All constructions precede \(L,w\).  Taking the original \(K\)
in CAF3 instead of the canonically isomorphic \(\mathbb Q(\delta)\)
preserves its original embedding and all the field conditions.
This claim verifies the seed-relevant part of \texttt{gen\_gd}; it
does not claim the unused full \texttt{Gen} fields as new results.
% END CAF9

% BEGIN CAF10
\paragraph{CAF10. The separate CountV host has its own seed modulus.}
\label{claim:CAF10}
For the original \(\operatorname{HostData}(\ell,e,\psi,o_1)\),
take
\[
 \eta=\eta_4,\qquad \Psi=\epsilon^e\psi,
 \qquad M=\ell N,
 \qquad Q=N_{\ell'}^{\,2}.
 \tag{CAF.14}
\]
For every complex conjugation \(\sigma\) and every \(B\) over
\(\ell\), these data satisfy \(\mathcal S\).  Thus CAF7 supplies
\texttt{SeedClause} for every compatible later \(L,w\), with its
original host characters.  This assertion includes \(\ell=3\).

\emph{Proof.}
The source \texttt{HostData} has \(\ell\ne2\), \(e,o_1\ge1\),
even \(N>0\), the parity \(\psi(-1)=(-1)^e\), and dyadic
conductor exponent \(c_2(\psi)\ge4\).  Its additional order and
size inequalities are retained but are unnecessary for this seed
construction.  The character \(\eta_4\) is the character of the
actual quadratic field \(\mathbb Q(i)\subset\Omega\); choose
\(\delta=i\), so \(D=1\).  It is unramified at odd \(\ell\)
and factors through \(\mu_4\).
Both \(\eta_4\) and \(\alpha\) factor through
\(\mu_{M^2}\), because \(4\mid M^2\), \(N\mid M^2\), and
\(\ell\mid M^2\).  Take \(a=2\).
The residue field is finite because \(O\) is finite over
\(\mathbb Z_\ell\), and the stated parity gives
\(\alpha(\sigma)=1\).  On inertia at \(\ell\), the residual
character of \(\psi\) kills the \(\ell\)-Sylow subgroup of
\((\mathbb Z/\ell^r\mathbb Z)^\times\), because \(k^\times\)
has order prime to \(\ell\).  It is therefore a power
\(\omega^b\), with one \(0\le b<\ell-1\), and the exponent
\(j=e+b\) works at every prolongation by conjugacy.  Thus no
prime-to-\(\ell\) order assumption on \(\psi\) is required.

Here is the needed noncontainment without using the
\(GD.\mathrm{gen}\) guard.  Write \(\psi_2\) for the dyadic
CRT component.  Its restriction to the subgroup of units congruent
to one modulo four is nontrivial; otherwise its conductor would
divide four, contradicting \(c_2\ge4\).  The image of that subgroup
has nontrivial power-of-two order.  Since the coefficient field has
characteristic zero, its unique element of order two is \(-1\).
Hence there is a dyadic unit \(v\equiv1\pmod4\) with
\(\psi_2(v)=-1\).  If \(\psi_2(-1)=1\), take \(u=-1\);
otherwise take \(u=-v\).  In both cases
\[
                  u\equiv-1\pmod4,\qquad \psi_2(u)=1.
 \tag{CAF.15}
\]
Total ramification of the dyadic cyclotomic extension supplies an
element \(s\) of inertia at two with this dyadic cyclotomic action.
All prime-to-two roots of unity are unramified there.  Thus the other
CRT components of \(\psi\) are one at \(s\), and the
\(\ell\)-adic cyclotomic character is one because \(\ell\) is odd.
Equation (CAF.15) gives
\(\omega(s)=\alpha(s)=1\) and \(\eta_4(s)=-1\), proving
\(\mathcal S\).  CAF6 gives
\(Q=(M^2\ell)_{\ell'}=N_{\ell'}^2\).
This verifies the seed-relevant part of \texttt{gen\_host} and
preserves the complete host small-prime scope.  It does not replace
the host by the \texttt{Dat} field or assume
\(\ell-1\nmid2e\).
% END CAF10

% BEGIN CAF11
\paragraph{CAF11. Installing the seed on the same marked frame.}
\label{claim:CAF11}
Choose \(Q\) by CAF6, CAF9, or CAF10 before selecting the frame.
Suppose the original \texttt{HasFrameW} at that \(Q\) has been
supplied.  Its very same
\(L,g_1,\Phi,C_0,w\), with the same \texttt{GalLink}, then
satisfies the full original \texttt{HasFrame}, including its
\texttt{SeedClause}.  In particular, an eventual proof of
\texttt{frame\_ex} can be consumed in both the \texttt{GdFrameEx}
and \texttt{HostFrameEx} applications without a new seed choice or
a new marked element.

\emph{Proof.}
Unpack the supplied \texttt{HasFrameW}.  It contains a
\texttt{Linked} witness, hence a literal
\(k:\operatorname{GalLink}(\eta,\Psi,\sigma,g_1,B,M,L,\Phi)\),
the existing \(\Phi.\mathrm{Ok}(C_0)\) and
\(\Phi.\mathrm{Big}(C_0)\), a witness prime \(w\),
\(w\equiv1\pmod Q\), and \(U_L\subseteq H_\ell(w)\).
The \texttt{ab} and \texttt{gal} fields of \(k\) give abelian
\(L\) and \(G_L=U_L\cap\ker\eta\).

The \texttt{g0} field gives \(g_1\in U_L\) with
\(p(g_1)=\Phi.g_0\).  The \texttt{del} field and the fact
\(\Phi.g_0\in\ker\Phi.\mathrm{del}\) imply
\(\eta(g_1)=1\).  Hence \(g_1\in G_L\).
The \texttt{psi} field identifies
\(\alpha(g_1)=\Phi.\Psi(\Phi.g_0)\), whose square is not one
by \(\Phi.\mathrm{Ok}.\mathrm{baseSubgroup\_spec}\).
Thus the final witness in (CAF.10) is the already selected \(g_1\).
CAF7 provides the full \texttt{SeedClause}.  Keep the same
\(L,g_1,\Phi,C_0,k\), its \texttt{Ok} and \texttt{Big} proofs,
and add that seed.  These are exactly all fields of
\texttt{HasFrame}; none has been reselected.

The source's \texttt{gd\_ex} and \texttt{host\_ex} use precisely
this installation after \texttt{frame\_ex}.  CAF9 and CAF10 give
their seed moduli directly, so their \texttt{Gen} constructions are
not additional seed hypotheses.  The existence of the large
\texttt{HasFrameW} is still required for either full frame-existence
conclusion.
% END CAF11

\subsection{The exact remaining interface}

The completed internal producers are \texttt{sc\_norm},
\texttt{sc\_teich}, \texttt{sc\_strip}, \texttt{sc\_quad},
\texttt{sc\_disc}, \texttt{hth}, \texttt{sc\_theta},
\texttt{auxw\_ex}, \texttt{sc\_cong}, \texttt{seed\_cl}, and
\texttt{field\_ex}, together with the two actual arithmetic
specializations of the seed hypotheses and the same-frame
installation.  The original \texttt{Gen} and \texttt{UFacts}
antecedents of the first two major source endpoints have been
preserved or weakened by proved implications; they have not been
silently assumed to exist.

For the remaining large frame, the source next chooses a real
quadratic subgroup \(H_0\), sets
\[
 U_m=H_0\cap H_\ell(w)\cap H_b\cap
          \ker\bigl(\epsilon\bmod\ell^{m+1}\bigr)^{\ell-1},
 \tag{CAF.16}
\]
proves its full \texttt{UFacts}, obtains one uniform class-rank
bound \(C_0\), and takes \(m=3C_0+7\).  Then
\(\ell^m>3C_0+7\).  CAF8 identifies the actual field belonging
to that very \(U_m\).  The still required pieces are the strict
cocycle constructions, the actual quotient/local linking frame,
and the cohomological \texttt{Ok} fields, with their exact
\texttt{GC.Hypotheses\_2} input.  The elementary inequality alone
does not prove those fields.

The CountV lane must keep \(M=\ell N\), \(\eta=\eta_4\), and
\(H_b=\ker(\kappa_B)\), where \(\kappa_B\) is the away-from-\(\ell\)
part of the \(\ell\)-power-order character.  The \texttt{Dat}
lane uses \(M=\ell N|\operatorname{disc}K|\) and \(H_b=G\).
Moreover, its auxiliary \(\Phi\) is the cohomological frame over
\(G_{L^+}\).  It is not the original CountV host frame over
\(G_{\mathbb Q,S}\).  Transport between them must retain their
actual projection maps, coefficient modules, local conditions, and
normalized Selmer lengths.  No \texttt{HOST\_FRAME\_LENGTH},
\texttt{HOST\_VAN\_E}, or \texttt{HOST\_XPP} closure is claimed here.

Finally, \texttt{LinkedEndEAlt} consumes both rank-two and rank-three
automorphic seeds, the general \texttt{ModUAlt} supplier, the same
frame and its common coefficient extension.  CAF7 supplies the
character triple those seeds require.  It does not assert that the
automorphic representations have been constructed.  Returning from
\texttt{GdEndAlt} to the full original \texttt{GdEnd} likewise
retains its local and descent matching.  Thus the selected whole
obligations \texttt{AUT\_UNRAM\_COEF}, \texttt{THM\_S6\_S}, and
the CountV joins remain open, while their common character-seed
producer is now completely written.
